\section{The odd-product quartics in 523 and 540 variables}\label{sec:v9}

The following construction incorporates the November 7 contribution of
Gabriel Conant, developed with ChatGPT. We give its arithmetic argument,
including the connection to the results proved above. Variable counts
include the input. The older quartics remain as source and proof baselines.

\begin{theorem}[The Version 9 quartics]\label{thm:v9-quartics}
There are specified integer polynomials $P_{523}$ and $P_{540}$, of total
degree four in respectively 523 and 540 variables, such that
\[
 q\in\Z\ \Longleftrightarrow\ \exists v\in\Q^{522}:P_{523}(q,v)=0
 \ \Longleftrightarrow\ \exists w\in\Q^{539}:P_{540}(q,w)=0.
\]
The first has 35,163 nonzero monomials and largest coefficient magnitude
$64\cdot3^{2^{27}}$. Its coefficient table uses the fixed integer
$S_0=3^{2^{26}}$. The second has 35,213 nonzero monomials and largest
coefficient magnitude $3^{2048}+1005072314631333<2^{3247}$.
Its complete decimal table occupies 708,848 bytes. Both have content one.
\end{theorem}

The proof is the composition of the source equivalences, local and global
arguments, and exact polynomial certificates below. The intermediate
554-variable construction is included to separate the odd-product
argument from the additional formal-depth refinement. It is not a further
claimed improvement over the two polynomials in the theorem.

\subsection{Inherited inputs and the exact changes of convention}\label{subsec:v9-inputs}

Use the normalized curve and generator of
Section~\ref{subsec:v5-point-source}, with $P=P_*$, not a new generator.
The coefficient basis is $(1,\sqrt2)$ and conjugation sends
$a+b\sqrt2$ to $a-b\sqrt2$. At unramified finite places $v(p)=1$;
at ramified places the integral valuation has value group $\Z$.
Heights are absolute logarithmic Weil heights, with the local degree
weights divided by $[K:\Q]=2$. Throughout this section set
$\mathcal H=H_{19}$ of Theorem~\ref{thm:quartic-height19}.

The source is the system in Section~\ref{subsec:v5-complete-source},
with the following explicit changes. Replace the comparison targets
$[4]P_*,[4]Q_i$ by $P_*,Q_i$; replace the orders 16 and 52 by 12 and 40
before the further refinement below. Retain the sixteen-test product,
the same five contacts, the scalar norm tree with orders $(1,1,1,1,32,32)$,
the Boolean correction and the preparation-zero pole equations.
Change $\beta$ to $2X^2+3$; use the point $S-P$ only in the scalar and
physical-anchor blocks; and change the fixed coefficient $\mathsf S$
to $S_0=3^{2^{26}}$. The point-chart replacement, the three new scales,
and the final order-one/order-forty test list are given below. All
unchanged quotient guards and permitted zero values of pole witnesses
are retained. This description, the displayed equations below, and the
pinned rational straight-line program specify the source without a
choice of arithmetic witnesses.

\begin{center}\small
\begin{tabular}{p{.29\textwidth}p{.63\textwidth}}
\toprule
Input & Exact result used and translation\\\midrule
Elliptic labels and domains & Section~\ref{subsec:v5-point-source},
Lemma~\ref{lem:v5-point-domains}, and \eqref{eq:o8-indexheight}:
$2E(K)=\Z P_*$, nonzero coordinate domains, and the same lower bound
$h(x([n]P_*))\ge n^2-6$.\\
Height of rational contacts & Theorem~\ref{thm:quartic-height19},
with $\mathcal H<2^{2093966}$, applied through the fractional linear
map and contact calculation of Lemma~\ref{lem:v5-height-transfer}.
Only that calculation is reused; its exponent 24 is replaced by 19.\\
Pole obstruction and realization & Lemmas~\ref{hyp:o8-pole},
\ref{lem:original-seven-simultaneous-realization}, and
\ref{lem:v5-one-place}. These include the restricted-converse dependency
proved in the earlier arithmetic sections; it has not been discarded.\\
Scalar and anchor & Lemma~\ref{v5:scalar32-bridge}, with its proof
applied to $2X^2+3$ and $S-P$, and
Lemma~\ref{lem:v5-three-witness-anchor} with the smaller exponent
calculated just below.\\
Local comparisons & Lemma~\ref{lem:v5-unramified-log} and the
recurrence of Lemma~\ref{lem:quartic-reciprocal-log}; orders and
target changes are justified below, not inferred by renaming indices.\\
Global choices & Theorem~\ref{thm:v5-simultaneous-primes},
Lemma~\ref{lem:v5-two-reductions}, Proposition~\ref{prop:v5-artin-patterns},
Section~\ref{subsec:v5-diagonal}, and
Lemma~\ref{v5:real-progression-bridge}. All fixed data precede the
asymptotic prime choice.\\
Final contradictions & Lemma~\ref{lem:o8-lucas} and the explicit
new height inequalities below.\\\bottomrule
\end{tabular}
\end{center}

The published results behind these inputs retain the versions and
convention qualifications in the theorem-application inventory. The new
contribution does not supply an assumed replacement for the restricted
converse, the height theorem, or the Gaussian prime argument.

\begin{lemma}[Direct targets and smaller anchor]\label{lem:v9-bridge}
For the changed source above, the local comparison and completeness
arguments remain valid with direct targets. The scalar anchor on
$R=[a-1]P_*$ implies
\[
 0\ne a-1\in2\,3^{2^{24}-1}\Z,
 \qquad \log_2|a-1|>25165823.5,
 \qquad h(\beta)\le\mathcal H|a-1|^{19/32}.
\]
It permits arbitrarily divisible positive scalar indices $a-1$.
\end{lemma}
\begin{proof}
For direct targets set $V_0=H-P_*$ and $V_i=U_i-Q_i$, with
$H=[h]P_*$, $U_i=[u_i]P_*$ and $Q_i=[n_i]P_*$. The quotient guards
give finite nonidentity points. At a consumed prime the full logarithm
ratios are $u_i/h$ and $(u_i-n_i)/(h-1)$. One of $h,h-1$ is a unit.
As in Lemma~\ref{lem:v5-two-anchors}, the cross congruence first makes
both approximate ratios integral, and then makes their relative
errors absolute. Subtraction of the two integral congruences gives
$L_i\equiv n_i$, without division by four. The logarithm recurrence
and denominator bound in Lemma~\ref{lem:quartic-reciprocal-log} give
the same estimates at orders 12 and 40. The test list below strengthens
the input order to 40 in the final refinement.

For completeness, the reduced elliptic groups have all four rational
two-torsion points, since $p>76176$ and the three roots $0,\pm d$
are distinct. Thus their doubling images have order $d_\pm(p)/4$.
The point $P_*$ is itself a double, so these integers annihilate its
reductions. The prime theorem and Lemma~\ref{lem:v5-two-reductions}
therefore give, on setting $M_\pm=\operatorname{lcm}_p(d_\pm(p)/4)$,
\[
 \gcd(M_-,M_+)=1,\quad\gcd(M_-,L)=16,\quad\gcd(M_+,L)=1.
\]
Here $64\mid L$, the exact two-parts are $2^6$ and $2^2$ before
division by four, and all odd prime divisors of $L$ were avoided.
Consequently $h\equiv0\pmod{M_-}$, $h\equiv1\pmod{M_+}$ and
$h\equiv h_0\pmod L$ are compatible when $64\mid h_0$.
The real-progression argument supplies large positive choices and
avoids the finite forbidden point and quotient loci. The later
prime-separation lemma justifies adding the factor $\Pi^{19}$.

For the anchor, repeat the proof of
Lemma~\ref{lem:v5-three-witness-anchor} with $C=3^{2^{25}}$ and
$\mathsf S=C^2=S_0$. Its quaternion criterion and anisotropic
binary form are unchanged. Thus the two rational coefficients of
$f(R)/C$ belong to $\Z_{(2)}\cap\Z_{(3)}$. The exact local calculation
of Section~\ref{subsec:o8-anchor-subgroup} now gives
$a-1\in2\,3^{2^{24}-1}\Z\setminus\{0\}$, and
\[
 \log_2|a-1|\ge1+(2^{24}-1)\log_2 3
 >1+\tfrac32(2^{24}-1)=25165823.5.
\]
For the scalar proof the new label is at least three, its five-contact
product never vanishes, and all its negative valuations at the
consumed odd primes are even. These are exactly the properties of
$X^2+2$ used in Lemma~\ref{v5:scalar32-bridge}; they hold for
$2X^2+3$ by the calculation below. That proof therefore gives
$M_5(\beta)^{32}\mid a-1$. Applying
Theorem~\ref{thm:quartic-height19} gives the asserted exponent $19/32$.
Its finite correction and continuity arguments apply to the same
norm tree with the scalar point $R$; choosing $a-1$ sufficiently
divisible meets the anchor and any further prescribed divisibility.
\end{proof}


\subsection{An exact replacement of the point chart}
Put $K=\Q(\sqrt2)$, $\delta=3-\sqrt2$ and $D=\delta^2=11-6\sqrt2$. The old chart, with letters in $K$, is
\begin{equation}\label{v9chart:eq:old}
 u=t^2,\qquad x=h^2,\qquad u(x-u/4)=D,\qquad y=h(u/2-x).
\end{equation}
The new chart is
\begin{equation}\label{v9chart:eq:new}
 x=h^2,\qquad j^2=x^2-D,\qquad y=hj.
\end{equation}
An equation over $K$ means its two rational component equations. The last expression in \eqref{v9chart:eq:new} need not be assigned its own coordinates.

\begin{lemma}\label{v9chart:lem:squares}
If $U,V\in K$, $U\ne0$, and $V^2=U^3+4DU$, then $U$ is a square in $K$.
\end{lemma}
\begin{proof}
The ring $\mathcal O=\Z[\sqrt2]$ is a PID: rounding both coefficients of an element of $K$ to nearest integers leaves a remainder of absolute norm at most $1/2$, giving a Euclidean algorithm. Its units are $\pm(1+\sqrt2)^n$, and every totally positive unit is a square. For completeness, multiply a positive unit by a power of $1+\sqrt2$ to place its first real embedding in $[1,1+\sqrt2)$. Its other embedding is the positive or negative reciprocal. Writing the unit as $a+b\sqrt2$ then forces $a,b\in\{0,1\}$, and the norm condition and excluded upper endpoint leave only the unit~1. The two signs of the fundamental unit under the real embeddings give the assertion about totally positive units.

Let $m=v_{\mathfrak p}(U)$, with integral normalization. Away from $(\sqrt2)$ and $(\delta)$ the two terms $U^3$ and $4DU$ have valuations $3m$ and $m$. If $m$ is odd, their unequal minimum is odd, a contradiction. At $(\sqrt2)$ these valuations are $3m$ and $m+4$, with equality only at the even value $m=2$. Again $m$ cannot be odd. At $(\delta)$ they are $3m$ and $m+2$, with equality only at $m=1$. In that exceptional case write $U=\delta u$ with $u$ a local unit. Then
\[
 V^2=\delta^3u(u^2+4).
\]
The residue field is $\mathbf F_7$, in which $-4$ is a nonsquare. Thus the right side has valuation three, also impossible. Every finite valuation of $U$ is even.

Under each real embedding, $D>0$ and $U(U^2+4D)=V^2$, so $U>0$. By the PID property, $U=\gamma^2\eta$ for $\gamma\in K^*$ and a unit $\eta$. This unit is totally positive, hence a square.
\end{proof}

\begin{proposition}\label{v9chart:prop:chart}
The charts \eqref{v9chart:eq:old} and \eqref{v9chart:eq:new} have exactly the same projected pairs $(x,y)\in K^2$. Every solution of the new chart has $h,j,x,y\ne0$, and $x$ is positive under both real embeddings.
\end{proposition}
\begin{proof}
Given \eqref{v9chart:eq:old}, put $j=t^2/2-h^2$. Then
\[
 j^2=h^4-t^2h^2+t^4/4=h^4-D,
\]
and $y=hj$. Conversely, given \eqref{v9chart:eq:new}, set $U=2(h^2+j)$ and $V=2hU$. The identity $j^2=h^4-D$ gives
\[
 U^2-4h^2U+4D=0,
\]
hence $V^2=U^3+4DU$. Also $U\ne0$, since $U=0$ would imply $D=0$. By Lemma~\ref{v9chart:lem:squares}, choose $t\in K$ with $t^2=U$. It follows that
\[
 t^2(h^2-t^2/4)=D,\qquad h(t^2/2-h^2)=hj,
\]
as required.

If $h=0$, then $j^2=-D$, contradicting either real embedding. If $j=0$, then $h^4=D$, so $h^2=\pm\delta$. Taking norms yields $\Norm(h)^2=7$, impossible for $\Norm(h)\in\Q$. The remaining assertions follow immediately.
\end{proof}

The proof is specific to this field and this value of $D$. It does not use the generally false assertion that square abscissa characterizes doubles on every elliptic curve. It also avoids adding a Mordell--Weil rank or generator hypothesis to the chart replacement.

\subsection{Using the chart without ordinate variables}
All nine point modules use the new chart. It has six rational coordinates: the two components of each of $h,j,x$, subject to four quadratic component equations. The old module had ten coordinates, including the components of $t,h,u,x,y$. Thus the change saves four coordinates per point, provided each downstream ordinate is consumed as the expression $hj$.

For the difference of chart points $Q,R$, retain the existing inverse equations for $x_Q-x_R$ and define the slope by
\[
 (x_Q-x_R)m=h_Qj_Q+h_Rj_R.
\]
Its numerator and defining equation are quadratic in the available coordinates. Define the difference abscissa by the existing square of $m$, and its ordinate by
\[
 y_{Q-R}=m(x_Q-x_{Q-R})-h_Qj_Q.
\]
This again is a quadratic assignment. The retained inverse ensures that the slope has exactly its old value.

For a chart point,
\[
 z=-x/y=-h/j.
\]
Thus the local parameter is defined by $jz=-h$. When a logarithm value is combined with its parameter, replace $y\ell=-xG$ by $j\ell=-hG$, where $G$ is the unchanged truncated logarithm polynomial. At the scalar ratio replace
\[
 x_H\rho=-y_Hz_S\qquad\text{by}\qquad h_H\rho=-j_Hz_S.
\]
All cancellations are reversible by Proposition~\ref{v9chart:prop:chart}. Difference-point denominators and all other guards remain unchanged. The fixed target has
\[
 h=-1+\tfrac74\sqrt2,\qquad j=\tfrac{39}{8}-\tfrac92\sqrt2;
\]
these give exactly its previous abscissa and ordinate.

Together with Proposition~\ref{v9chart:prop:chart}, these identities give rational witness maps in both directions. To reconstruct the old chart, choose the square root supplied by Lemma~\ref{v9chart:lem:squares}; then evaluate all old assignments. The choice of sign does not alter $(x,y)$ or any downstream value.

\subsection{Shorter weight circuits}\label{subsec:v9-laurent}
For each of nine detectors, write $e$ for its nonzero weight and $q$ for its order-one weight. The required outputs are $q_{12}=qe^{22}$ and $q_{40}=qe^{78}$. The old construction uses seven pure powers of $e$ and two final quotients. Instead construct
\[
 a_1=e^{-1},\quad a_3=a_1^2/e,\quad a_7=a_3^2/e,\quad
 a_{14}=a_7^2,\quad a_{29}=a_{14}^2/e,\quad a_{55}=a_{29}^2/a_3,
\]
and impose
\[
 a_{29}q_{12}=qa_7,\qquad a_{55}q_{40}=q_{12}e.
\]
These are six power coordinates and two output coordinates. They give the same outputs and save one coordinate per detector.

The nonzero condition is inherited, rather than presumed from a quotient notation. In the source, $w_0$ has a retained inverse and
\[
 w_0n=(w_0^2-2w_1^2)/2,\qquad (e\Delta)b=n^2.
\]
Since $w_0\ne0$ and 2 is a nonsquare in $\Q$, the first equation gives $n\ne0$. The second gives $e\Delta\ne0$. Therefore every Laurent denominator above is nonzero. Conversely all Laurent coordinates are uniquely defined for each source solution. The changes introduce no new zero branches and exclude no source solution.

\subsection{Scalar preparation with one fewer coordinate}\label{subsec:v9-preparation}
The scalar pole uses a nonzero input $b$ and a preparation witness $t$, with
\[
 U=\frac{b(1-bt^2)}{1+bt^2}.
\]
The old circuit stores $b^{-1},t_2=t^2,bt_2,U$. Replace these four coordinates by $t_2,T,V$, with
\[
 t_2=t^2,\qquad b(T-t_2)=1,\qquad TV=2,
\]
and use the affine expression $U=V-b$. The equation $b(T-t_2)=1$ forces $b\ne0$, and $TV=2$ forces $T\ne0$. Consequently $T=t^2+b^{-1}$, $1+bt^2=bT\ne0$, and
\[
 U=2/T-b=\frac{b(1-bt^2)}{1+bt^2}.
\]
Conversely every old solution extends uniquely by these values of $T,V$. Thus the modification saves one coordinate and preserves all downstream quantities. It also supplies the nonzero condition formerly imposed by the separate coordinate $b^{-1}$.


\subsection{An exact product-height identity}\label{v9:sec:height}
Put
\[
 K=\Q(\sqrt2),\quad d=3-\sqrt2,\quad D=d^2=11-6\sqrt2,
 \qquad E:y^2=x^3-Dx,
\]
and put
\[
 R_0=(1+\sqrt2,2\sqrt2),\qquad
 P=2R_0=\left(\frac{57-28\sqrt2}{8},\frac{-660+417\sqrt2}{32}\right).
\]
The model calculation in Section~\ref{sec:degree8-original} proves that $P=P_*$ has infinite order and $2E(K)=\Z P_*$; this is the same normalized curve and generator, not the earlier model of Proposition~\ref{prop:generator}. Write $\operatorname{h}$ for the absolute logarithmic Weil height. The following lemma uses the displayed coordinates and doubling identity, together with the infinite order of $P$.

\begin{lemma}\label{v9:lem:product}
For odd integers $n_1,\ldots,n_r$,
\[
 \operatorname{h}\!\left(\prod_{i=1}^r x([n_i]P)\right)
 =\sum_{i=1}^r\operatorname{h}(x([n_i]P)).
\]
\end{lemma}
\begin{proof}
Let $T_0=(0,0)$. At a finite place $v$ with $v(D)=0$, use the finite quotient $E(K_v)/E_1(K_v)$, where $E_1$ denotes the formal neighborhood of infinity for this integral equation.
Here we use the integral formal group construction of
\cite[Chapter IV, Sections 1--3]{silverman}. Its parameter is $z=-x/y$,
and $x=z^{-2}(1+O(z))$, $y=-z^{-3}(1+O(z))$ identify this neighborhood
with the maximal ideal. It is an open subgroup; the proper curve
$E(K_v)$ is compact, so its quotient is finite. This argument does not
assume that the equation is minimal at the place above two. For a finite nonidentity point $Q$, membership in $E_1$ is equivalent to $v(x(Q))<0$. Since
\[
 x(Q+T_0)=-D/x(Q),
\]
positive valuation of $x(Q)$ is equivalent to $Q+T_0\in E_1$. The image of $T_0$ in this quotient has order two.

If an odd multiple of $P$ belongs to $E_1$, the image of $P$ has odd order. None of its multiples can then equal the image of $T_0$, so all $x([n_i]P)$ have nonpositive valuation. If no odd multiple belongs to $E_1$, all these abscissae have nonnegative valuation. Thus the valuations have a common weak sign. This argument includes the place above two.

The only place with $v(D)>0$ is $\mathfrak d=(3-\sqrt2)$ above seven. The displayed integral model is minimal there, since $v_{\mathfrak d}(-D)=2<4$. Its reduction is $y^2=x^3$ over $\F_7$. Substituting $\sqrt2=3$ shows that $P$ reduces to $(1,6)$, a nonsingular point. By \cite[Proposition VII.2.1, pp. 188--190]{silverman}, all multiples therefore belong to the subgroup $E_0(K_{\mathfrak d})$ of points with nonsingular reduction. A positive valuation of $x$ would force a positive valuation of $y$, and hence reduction to the singular point $(0,0)$. It cannot occur. The valuations again have a common weak sign.

At either real embedding, each $[n_i]P$ is a double of a non-torsion point. The duplication formula makes its abscissa a nonzero square. Its ordinate is nonzero, so the equation of $E$ gives $x([n_i]P)>d$ at that embedding. Both real values $3\pm\sqrt2$ of $d$ exceed one.

At a finite place, numbers whose valuations have the same weak sign satisfy
\[
 \log\max\left(1,\left|\prod_i x([n_i]P)\right|_v\right)
 =\sum_i\log\max(1,|x([n_i]P)|_v).
\]
At both real places the same equality holds because all the factors exceed one. Summing with normalized local weights proves the assertion.
\end{proof}

The product here is an element of $K$. It is not its norm down to $\Q$. Norms can cancel a pole at one split place against a zero at the other, whereas the proof above works separately at every place.

\begin{corollary}\label{v9:cor:scales}
Suppose the four root points $Q_i=[n_i]P$ have odd indices, and write
\[
 C=\prod_{i=1}^4x(Q_i)=C_0+C_1\sqrt2.
\]
The original nine height-detector scales may be replaced by
\[
 1,\qquad C_0^2,\qquad (1+C_1^2)^2.
\]
The root-height estimate retains the original constant $7\mathcal H$, where $\mathcal H$ is the exponent-19 height constant of the source.
\end{corollary}
\begin{proof}
Both real values of $C$ are positive, so $C_0>0$. All three scales are positive nonzero rational squares. If a nonzero integer $N$ is divisible by the $k$th powers of the three contact radicals, the unchanged detector estimate gives
\[
 \operatorname{h}(\lambda),\ \operatorname{h}(C_0^2\lambda),\
 \operatorname{h}((1+C_1^2)^2\lambda)
 \le T:=(\mathcal H+\log3)|N|^{19/k}.
\]
The elementary rational height identities yield
\[
 \operatorname{h}(C_0)\le T,\qquad \operatorname{h}(C_1)\le T/2.
\]
Indeed, height subadditivity first bounds the heights of the two scales by $2T$; also $\operatorname{h}(C_0^2)=2\operatorname{h}(C_0)$ and $\operatorname{h}((1+C_1^2)^2)\ge4\operatorname{h}(C_1)$. Hence
\[
 \operatorname{h}(C)\le\operatorname{h}(C_0)+\operatorname{h}(C_1)+\tfrac32\log2
 \le\tfrac32T+\tfrac32\log2
 <7\mathcal H |N|^{19/k}.
\]
Here $\mathcal H\ge1$. Lemma~\ref{v9:lem:product} bounds each root height by $\operatorname{h}(C)$. In particular, the inherited lower bound, with $B=\max_i|n_i|$, gives
\begin{equation}\label{v9:eq:height-transfer}
 B^2-6<7\mathcal H |N|^{19/k}.
\end{equation}
The old completeness proof uses its scales only after the roots have been fixed, and requires fixed positive rational squares. Thus its prime selection and pole construction apply to this shorter list.
\end{proof}

\subsection{Enforcing odd indices by quadratic equations}\label{v9:sec:chart}
Put $\epsilon=1+\sqrt2$. The existing chart for a point $Q$ has $K$-coordinates
\[
 x=h^2,\qquad j^2=x^2-D,\qquad Q=(x,hj).
\]
To this chart add $t,r\in K$ and the equations
\begin{equation}\label{v9:eq:odd-chart}
 t^2=2(x+j),\qquad th+x+j=\epsilon r^2.
\end{equation}
These are four rational quadratic equations in four additional rational coordinates. No extra inverse witnesses are needed.

\begin{proposition}\label{v9:prop:odd-chart}
The projected points of the augmented chart are exactly $P+4E(K)$, that is, the odd multiples of $P$.
\end{proposition}
\begin{proof}
Use the squareclass homomorphism
\[
 \alpha:E(K)\longrightarrow K^*/K^{*2},\qquad
 \alpha(O)=1,\quad \alpha(T_0)=[-D]=[-1],\quad
 \alpha((a,b))=[a]\quad(a\ne0).
\]
For clarity, its homomorphism property follows from the line-intersection identity
\[
 x(R+S)x(R)x(S)=
 \left(\frac{x(R)y(S)-x(S)y(R)}{x(R)-x(S)}\right)^2
\]
when the displayed abscissae are nonzero and distinct. Doubling follows from
$x(2R)=((x(R)^2+D)/(2y(R)))^2$; addition by $T_0$ follows from $x(R+T_0)=-D/x(R)$. Inverse pairs and identity summands are immediate. If $R+S=T_0$, then $x(R)x(S)=-D$. These cases cover the exceptional pairs.

Since $2E(K)=\Z P$ and $2R_0=P$, every point of $E(K)$ has a unique expression $mR_0+T$, where $m\in\Z$ and $T\in E[2](K)$. The two-torsion is generated by $T_0$ and $T_d=(d,0)$, so
\[
 \alpha(mR_0+aT_0+bT_d)=[\epsilon]^m[-1]^a[d]^b.
\]
The classes $[\epsilon],[-1],[d]$ are independent. A square $\epsilon^c(-1)^a d^b$, with $a,b,c\in\{0,1\}$, would have rational square norm $(-1)^c7^b$. This forces $c=b=0$, and the real embeddings then force $a=0$. Consequently
\begin{equation}\label{v9:eq:squareclass}
 \alpha(R)=[\epsilon]\quad\Longleftrightarrow\quad R\in R_0+2E(K).
\end{equation}

Given the augmented chart, set $u=t^2=2(x+j)$ and
\[
 a=th+u/2,\qquad b=at.
\]
Expansion gives
\[
 u(x-u/4)=D,\qquad a^2-ua=D,\qquad b^2=a^3-Da.
\]
In particular $t,a,b$ are nonzero. The tangent slope at $R=(a,b)$ is $h+t$; duplication gives $2R=(x,hj)=Q$. Equation~\eqref{v9:eq:odd-chart} says $a=\epsilon r^2$, so $\alpha(R)=[\epsilon]$. By~\eqref{v9:eq:squareclass}, $Q\in P+4E(K)$.

Conversely, for $Q=(2m+1)P$, take $R=(2m+1)R_0=(a,b)$. Then $a,b\ne0$, and $a=\epsilon r^2$ for some $r\in K$. Set
\[
 t=b/a,\qquad h=(a^2+D)/(2b),\qquad x=h^2,
 \qquad j=t^2/2-x.
\]
The same duplication identities give $Q=(x,hj)$, $j^2=x^2-D$, and both equations~\eqref{v9:eq:odd-chart}.
\end{proof}

Only the four root points receive this augmentation. The shared point and the four comparison points retain their old charts: their indices need not be odd. The total cost is sixteen coordinates and sixteen quadratic residuals. The projected-point formulation matters, since the added equations can select a sign representation of the old $h,j$ chart.

\subsection{Odd labels, scalar bounds, and completeness}\label{v9:sec:labels}
Write the root points as
\[
 A=[n]P,\qquad S=[a]P,\qquad T=[t]P,\qquad Z=[b]P.
\]
Here $S$ denotes a point; the fixed coefficient remains $S_0$. Replace the old input expression by
\[
 \beta=2X^2+3.
\]
For rational $X$, integrality of $\beta$ is equivalent to integrality of $X$: if $X=u/v$ in lowest terms, then $\beta\in\Z$ forces $v^2\mid2$, and hence $v=1$. For integral $X$, $\beta$ is odd.

In the scalar norm block and the physical anchor, replace $S$ by
\[
 R=S-P=[a-1]P.
\]
The comparison and product-height blocks still use $S$. The retained subtraction guard excludes $a=\pm1$, which is harmless for completeness. The original scalar contact product is
\[
 \beta(\beta-1)(\beta-2401)(32\beta+343)(2\beta-1).
\]
It never vanishes for rational $X$: $\beta\ge3$, and $\beta=2401$ would give $X^2=1199$, which is not a rational square. At an odd prime, a negative valuation of $\beta$ is twice that of $X$, as required by the scalar argument.

\subsubsection{Soundness}\label{subsec:v9-soundness}
Put $c=|a-1|$. Applying the scalar and anchor lemmas to $R$ gives
\[
 \operatorname{h}(\beta)\le\mathcal H c^{19/32},\qquad
 \log_2c>25{,}165{,}823.5,
 \qquad 1\le\mathcal H<2^{2{,}093{,}966}.
\]
The anchor is still $S_0=3^{2^{26}}$. These explicit bounds imply
\[
 c>(80\mathcal H)^{10},\qquad
 |a|\ge c-1\ge c/2>(40\mathcal H)^{10},
\]
and
\begin{equation}\label{v9:eq:margins}
 \operatorname{h}(\beta)<c/80\le |a|/40,
 \qquad\log(40\mathcal H)<c/80\le |a|/40.
\end{equation}
For example $10(7+2{,}093{,}966)<25{,}165{,}823.5$, and $80<2^7$. The first estimate in~\eqref{v9:eq:margins} follows from $c^{13/32}>80\mathcal H$; the logarithmic estimate is weaker still.

The local comparison lemmas retain their original conclusions. Let
\[
 F=t^2-a(a-2)tb+(a-2)^2b^2-1,
 \qquad B=\max(|n|,|a|,|t|,|b|).
\]
For $B>4$, the triangle inequality gives $|F|\le2B^4+6B^3+5B^2+1<5B^4$. If $F\ne0$, the three radicals divide it to order 40, so~\eqref{v9:eq:height-transfer} gives
\[
 B^{1/10}<7\mathcal H\,5^{19/40}+6<24\mathcal H.
\]
This contradicts the anchor. Thus $F=0$. Since the point $Z$ is finite and nonidentity, $b\ne0$. Lemma~\ref{lem:o8-lucas} now gives $\log B>|a|-3$.

If $n\ne\beta$, the nonzero numerator $N$ of $n-\beta$, in reduced form, satisfies $|N|\le2e^{\operatorname{h}(\beta)}B$, and the radicals divide it to order 12. Equation~\eqref{v9:eq:height-transfer} and $B>4$ give
\[
 B^{5/12}<40\mathcal H e^{(19/12)\operatorname{h}(\beta)},
\]
so
\[
 \log B<\frac{12}{5}\log(40\mathcal H)+\frac{19}{5}\operatorname{h}(\beta)
 <\frac{31}{200}|a|<|a|-3.
\]
This is the other contradiction. Therefore $n=\beta\in\Z$, and $X\in\Z$.

\subsubsection{Completeness}\label{subsec:v9-completeness}
Fix $X\in\Z$, and set $n=\beta$. Perform the old scalar preparation: choose its finite correction and then a positive shared index $h_0$ divisible by 64 and sufficiently deep at the fixed small primes. Choose a positive scalar index $s$ sufficiently divisible to satisfy the scalar leaves and physical anchor for $[s]P$. We may require $6\mid s$ and take $s$ arbitrarily large. Set $a=s+1$. Then $a$ is odd, $a\equiv1\pmod3$, and $S-P=[s]P$ is exactly the prepared scalar point.

Define
\[
 U_0=0,\qquad U_1=1,\qquad U_{j+1}=aU_j-U_{j-1},
 \qquad k=2(a-2),
\]
and take $b=U_k/(a-2)$ and $t=U_{k+1}$. Modulo $a-2$, induction gives $U_j\equiv j$, so $b\in\Z$. Modulo two the sequence has period $0,1,1$; since $k\equiv1\pmod3$, both $U_k$ and $U_{k+1}$ are odd. The positive odd denominator $a-2$ shows that $b$ is odd. The Lucas identity gives
\[
 t^2-a(a-2)tb+(a-2)^2b^2=1.
\]
Thus all four root indices $n,a,t,b$ are positive and odd. Proposition~\ref{v9:prop:odd-chart} supplies their chart witnesses.

Fix these roots and the three scales of Corollary~\ref{v9:cor:scales}. Local continuity of the scalar and comparison functions at the prepared shared point gives the same fixed modulus $L$, with $64\mid L$. Only then choose the contact primes by the inherited simultaneous-prime theorem. The direct-target construction chooses a common index $h$ satisfying the primary and opposite reduction-order congruences and $h\equiv h_0\pmod L$, with the required real bounds. Its comparison points are $H=[h]P$ and $U_i=[hn_i]P$. Their shifted points have indices $h-1$ and $(h-1)n_i$.

All quotient-domain arguments remain valid. The product scales are fixed positive squares, so the original pole witnesses apply unchanged. The scalar block is preserved by the congruence modulo $L$. This gives a rational solution for every integral $X$ and completes the arithmetic proof of the primary construction.

\subsection{Replacing five logarithms by asymmetric formal depth}\label{v9:sec:depth}
This section is used only for the 523-variable construction. Let
\[
 U_0=H=[h]P,\quad U_i=[u_i]P,\quad
 V_0=H-P,\quad V_i=U_i-Q_i,
\]
where $Q_i=[n_i]P$ are the four roots. Put $f=1/x$ and $z=-x/y$. Keep the alternate truncated logarithms $\ell_{40}(V_i)=z(V_i)G_9(Df(V_i)^2)$, but replace every primary truncated logarithm by $z(U_i)$. Thus
\[
 L_i=z(U_i)/z(H),\qquad
 L_i'=\overline{\ell_{40}(V_i)/\ell_{40}(V_0)}.
\]
The bar denotes conjugation in $K$. The sixteen tests now have the following orders:
\[
\begin{array}{c|l}
1 & f(U_1),\ldots,f(U_4),\ \overline{f(V_0)},\ldots,\overline{f(V_4)}\\[2pt]
40 & f(H),\ L_1-L_1',\ldots,L_4-L_4',\ L_1-\beta,\
       \text{the direct Lucas residual}.
\end{array}
\]
The scalar tests retain their original order 32. The stronger beta tests imply all former order-12 conclusions.

\subsubsection{Local soundness}\label{subsec:v9-local-depth}
At a consumed unramified place above $p>76176$, a formal point $Q$ with $r=v(z(Q))\ge1$ satisfies
\begin{equation}\label{v9:eq:zerror}
 \log_E(Q)/z(Q)\in1+p^{4r}\mathcal O.
\end{equation}
This is the same formal-logarithm expansion used in the source, starting its tail at the first nonconstant term. Namely $v(f(Q))=2r$, and for its coefficients $a_j$ the factorial estimate is $v(a_j)\ge-(4j+1)/(p-1)$. If $4j+1<p$, the coefficient is integral. For the remaining $j\ge19044$, the displayed bound still makes the $j$th term have valuation at least $4r$. Both $z(Q)$ and its full logarithm have valuation $r$.

The new test on $f(H)$ gives $v(z(H))\ge20$. The order-one tests put the other points on the formal branch. The full logarithm ratios at the primary and conjugate places are
\[
 r_i=u_i/h,\qquad r_i'=(u_i-n_i)/(h-1).
\]
At least one of $h,h-1$ is a $p$-unit, so at least one of these ratios is integral. Even before knowing the depth of $U_i$, the weak unit estimate~\eqref{v9:eq:zerror} and the alternate truncation estimate show that $L_i,L_i'$ have the same valuations as $r_i,r_i'$. The order-40 cross test transfers integrality to both sides. Now $L_i$ integral implies $v(z(U_i))\ge v(z(H))\ge20$. The primary logarithmic corrections are therefore $1$ modulo $p^{80}$. The alternate corrections are $1$ modulo $p^{40}$. Since the ratios are integral, these relative estimates become absolute:
\[
 L_i\equiv r_i\equiv r_i'\pmod {p^{40}}.
\]
Multiplying by $h$ and $h-1$ and subtracting gives $L_i\equiv n_i\pmod {p^{40}}$. At an inert prime conjugation preserves the local valuation, so the same proof applies. The height contradictions of Section~\ref{v9:sec:labels} remain unchanged.

\subsubsection{A stronger choice within the existing prime-pattern argument}
For completeness, the common index must acquire extra primary divisibility. The following observation ensures compatibility with the opposite reduction orders.

\begin{lemma}\label{v9:lem:separation}
The contact primes in the source completeness argument may be chosen so that no selected prime divides any opposite reduced group order.
\end{lemma}
\begin{proof}
After merging repeated contact forms and fixing all scale and congruence data, the prime-pattern proof uses distinct forms $p_j=u+c_jw^2$, with fixed coefficients $c_j$. Integer dilations and ray congruences change only fixed constants. It counts $u\in[N/2,N]$ and $1\le w\le T$, where $T=\lfloor\tau(N)\sqrt N\rfloor$ and $\tau(N)\to0$ can be chosen arbitrarily slowly in the final diagonal argument. Its weighted count, after the prescribed obstruction removal, retains a fixed positive proportion of $NT$. These are the fixed-field weighted count in Proposition~\ref{prop:v5-artin-patterns} and the final diagonal argument of Section~\ref{subsec:v5-diagonal}, including its obstruction removal.

If there are $k$ forms, impose also
\[
 \frac{\tau(N)N^{1/6}}{(\log N)^k}\longrightarrow\infty.
\]
This is compatible with the same diagonal choice: for each fixed positive value of $\tau$, its threshold can be added to the existing finite list of thresholds. The strip $w<N^{1/3}$ has at most $O(N^{4/3})$ integer pairs and weight at most $O((\log N)^k)$ per pair. Its total weight is $o(NT)$, so an admissible pattern remains outside it.

For these patterns $w^2\ge N^{2/3}$ and $w^2=o(N)$. All $p_j$ are asymptotic to the same $u$. Hasse's bound \cite[Theorem V.1.1, p. 138]{silverman} gives $d_+(p_j)=p_j+O(\sqrt N)$ for the opposite reduced group order. Every ratio $d_+(p_j)/p_i$ is eventually strictly between $1/2$ and $2$. Hence $p_i\mid d_+(p_j)$ would force equality. If $i\ne j$, this contradicts
\[
 |p_i-p_j|\ge c_*N^{2/3},\qquad
 c_*:=\min_{i\ne j}|c_i-c_j|>0.
\]
If $i=j$, equality is impossible because $4\mid d_+(p_i)$ whereas $p_i$ is odd.
\end{proof}

\subsubsection{Completeness with deeper primary points}\label{subsec:v9-deep-completeness}
Keep the choice order in Section~\ref{v9:sec:labels}: the roots, scales and continuity modulus $L$ precede the prime selection. Choose primes using Lemma~\ref{v9:lem:separation}, beyond all prime factors of $L$ and the fixed scale data. Put
\[
 M_\pm=\mathop{\mathrm{lcm}}_p(d_\pm(p)/4),\qquad \Pi=\prod_p p.
\]
The direct-target prime theorem gives
\[
 \gcd(M_-,M_+)=1,\quad
 \gcd(M_-,L)=16,\quad\gcd(M_+,L)=1.
\]
Lemma~\ref{v9:lem:separation} gives $\gcd(\Pi,M_+)=1$, and the cutoff gives $\gcd(\Pi,L)=1$. Thus
\[
 h\equiv0\pmod{\Pi^{19}M_-},\qquad
 h\equiv1\pmod{M_+},\qquad h\equiv h_0\pmod L
\]
are compatible, since $64\mid h_0$. The inherited real-progression argument supplies a sufficiently large $h$ in this progression with the same real bounds.

The point $P$ is globally a double, so multiplication by the primary doubling-image order first makes it formal, with $z$-valuation at least one. Multiplication by $p^{19}$ raises this valuation by nineteen, using the formal logarithm isometry of \cite[Theorem IV.6.4, p. 132]{silverman}; these completions are unramified, $v(p)=1$, and $1>1/(p-1)$. Consequently $H=[h]P$ and $U_i=[hn_i]P$ have primary depth at least twenty, and their parameter ratios equal $n_i$ to order eighty. At the opposite places, $V_0=[h-1]P$ and $V_i=[(h-1)n_i]P$ are formal and their retained logarithm ratios equal $n_i$ to order forty.

The pole-completeness estimates still apply to these fixed test orders. In the source notation their large-real-parameter decay is $\lambda^{-3}$ for each order $m$, and the small-prime lower bound is $3\nu+(10m-8)\eta-2c$, which still tends to infinity with the preparatory divisibility. Scalar conditions and the anchor are preserved by $h\equiv h_0\pmod L$. This proves completeness for the optional refinement.

\subsection{Polynomial construction and exact verification}\label{subsec:v9-finite}
The generators apply three changes to the earlier point-chart source: the four odd-coset charts, the product and three scales, and the shifted scalar point with $\beta=2X^2+3$. The depth variant first changes the logarithms and test orders as in Section~\ref{v9:sec:depth}. Each field equation is represented by two rational component equations. Intermediate products are assigned coordinates as needed to keep every residual quadratic.

The full-log source has 577 raw coordinates; the pinned upstream compiler reduces this to 561. Seven further certified eliminations give 554. The depth source has 540 raw coordinates; compilation gives 532 and nine certified eliminations give 523. Every elimination is justified by exact polynomial identities and reversible substitutions preserving the input. Clearing residual denominators and taking the sum of their squares gives the displayed integer quartics. Over $\Q$, the sum vanishes exactly when every residual does.

The constant-chain versions add $c_0,\ldots,c_{16}$ with
\[
 c_0=3^{1024},\qquad c_{i+1}=c_i^2\quad(0\le i<16).
\]
These witnesses have unique rational values and $c_{16}=S_0$. Substitution for $S_0$ leaves every residual quadratic. Direct expansion gives the 3,247-bit maximum coefficient bound.

The accompanying archive contains the complete polynomial tables, residual systems, coordinate maps, elimination certificates, pinned source generators, arithmetic notes and independent checkers. Each final package has a standard-library-only \texttt{verify\_all.py}. It replays all eliminations, verifies the unique constant lift, independently expands the sums of squares, and compares every coefficient. Both full packages have passed fresh independent runs. A separate source replay rebuilds the two compiled precursors and checks the match to their frozen residual systems.

These finite checks certify the counts, degrees and stated algebraic transformations. The arithmetic assertions in Sections~\ref{v9:sec:height}--\ref{v9:sec:depth} received separate adversarial AI review; they have not received independent human peer review or a complete formal proof check. The inherited arithmetic inputs are the results of this manuscript identified in Section~\ref{subsec:v9-inputs}; no theorem from the contribution archive is assumed as an additional arithmetic input.



\begin{proof}[Proof of Theorem~\ref{thm:v9-quartics}]
The chart identities and Lemma~\ref{lem:v9-bridge} identify the exact
source used by the generators with the arithmetic system described here.
Proposition~\ref{v9:prop:odd-chart} enforces the odd root indices, so
Corollary~\ref{v9:cor:scales} supplies the required height bound.
Section~\ref{subsec:v9-local-depth} verifies the local implications for
the final test orders. Section~\ref{subsec:v9-soundness} then forces
the input to be an integer. Conversely, for an integer input,
Sections~\ref{subsec:v9-completeness} and
\ref{subsec:v9-deep-completeness}, using Lemma~\ref{v9:lem:separation},
construct witnesses for every source equation and retained guard.
The certified substitutions preserve this projection in both directions.
The sum of squares vanishes over $\Q$ precisely when all its residuals
vanish. Finally the 17-coordinate constant chain has exactly one rational
solution and ends at $S_0$, so it preserves the same projection. This
proves the statements for both polynomials.
\end{proof}

The files \nolinkurl{quartics/quartic-523-symbolic-coefficients.json}
and \nolinkurl{quartics/quartic-540-fully-expanded.txt} give the complete
tables. Their SHA-256 digests, respectively, are
\begin{center}\small
\nolinkurl{6a3685980de4512bdc7f3a781d0c0d4b32cce5829ec83f03f318264931bdbbe3}\\
\nolinkurl{83e7a31e11f9ad2edc2f67935e50bff437d87ef4b77d37b6a4391e8117770680}.
\end{center}
The reproduction command is
\nolinkurl{python scripts/check_odd_product.py}; the complete source,
certificates and repaired, non-mutating verifiers are in
\nolinkurl{quartics/quartic-523-540/}.

